% !TEX program = xelatex
\documentclass[12pt,a4paper]{article}
\usepackage{fontspec}
\usepackage{polyglossia}
\setmainlanguage{english}
\setmainfont{DejaVu Serif}
\setsansfont{DejaVu Sans}
\setmonofont{DejaVu Sans Mono}
\usepackage{geometry}
\geometry{margin=2.5cm}
\usepackage{amsmath,amssymb,amsthm,mathtools,bm}
\usepackage{microtype}
\usepackage{enumitem}
\usepackage{hyperref}
\hypersetup{colorlinks=true, linkcolor=blue, citecolor=blue, urlcolor=blue}
\newtheorem{definition}{Definition}
\newtheorem{axiom}{Axiom}
\newtheorem{theorem}{Theorem}
\newtheorem{proposition}{Proposition}
\newtheorem{remark}{Remark}
\DeclareMathOperator{\Adm}{Adm}
\DeclareMathOperator{\End}{End}
\DeclareMathOperator{\Tr}{Tr}
\DeclareMathOperator{\Attr}{Attr}
\DeclareMathOperator{\Struct}{Struct}
\DeclareMathOperator{\Dist}{Dist}
\DeclareMathOperator{\Valid}{Valid}
\DeclareMathOperator{\Auto}{Auto}
\DeclareMathOperator{\Input}{Input}
\DeclareMathOperator{\Dir}{Dir}
\DeclareMathOperator{\AutoAct}{AutoAct}
\DeclareMathOperator{\Coh}{Coh}
\DeclareMathOperator{\Sync}{Sync}
\DeclareMathOperator{\Loss}{Loss}
\DeclareMathOperator{\Stab}{Stab}
\newcommand{\M}{\mathcal M}
\newcommand{\B}{\mathcal B}
\newcommand{\Hcal}{\mathcal H}
\newcommand{\Ispace}{\mathfrak I}
\newcommand{\XiAuto}{\Xi_I^{\mathrm{auto}}}
\newcommand{\XiMan}{\mathcal M_{\Xi}^{\mathrm{auto}}}
\newcommand{\NXi}{\mathcal N_{\Xi}}
\newcommand{\PiAmp}{\Pi_{\mathrm{amp}}}
\newcommand{\PiStr}{\Pi_{\mathrm{str}}}
\newcommand{\PiStruct}{\Pi_{\mathrm{struct}}}

\title{\textbf{Automatically Structured Interface Fields in TMI-Physics 2.2}\\
\large From Quantum Possibility Distributions to Scalable Interface Manifolds}
\author{Salkutsan Aleksey}
\date{2026}

\begin{document}
\maketitle

\begin{abstract}
This paper formulates the physical layer of the Theory of Manifold Interfaces at the level of TMI-Physics 2.2. The minimal scalar interface field \(\chi_I\), introduced in earlier physical modules, is preserved as a low-dimensional amplitude projection of a fuller automatically structured interface field \(\Xi_I^{\mathrm{auto}}\). The central chain is: geometry defines admissible transitions; quantum theory distributes possible transitions; the primary interface provides the ontological ground of distinction; and an automatically structured interface field organizes distributed possibilities into stable structures. The paper distinguishes internal theorems of the model from physical hypotheses and experimental consequences.
\end{abstract}

\tableofcontents
\newpage

\section{Purpose and Status}
This document should be read as a physical module depending on TMI-Core 0.2, where automaticity is introduced as a definitional property of valid interfaces. The goal is not to prove that \(\Xi_I^{\mathrm{auto}}\) already exists in nature. The goal is to formulate the internal physical architecture that follows if a physical interface field is admitted inside TMI-Physics.

\section{Minimal Physical Chain}
The minimal earlier physical chain was
\begin{equation}
\boxed{g\Rightarrow\Adm_c(g)\Rightarrow\Hcal_I(g)\Rightarrow\Psi_I\Rightarrow\Dist_{\mathrm{poss}}(\Psi_I;g)\Rightarrow\chi_I\Rightarrow\Struct_I.}
\end{equation}
In TMI-Physics 2.2 this is not discarded. It is reinterpreted as a projection of the fuller chain
\begin{equation}
\boxed{g\Rightarrow\Adm_c(g)\Rightarrow\Hcal_I(g)\Rightarrow\Psi_I\Rightarrow\Dist_{\mathrm{poss}}(\Psi_I;g)\Rightarrow\XiAuto\Rightarrow\Struct_I.}
\end{equation}
The scalar field \(\chi_I\) remains the minimal observable amplitude projection of the full field.

\section{Interpretive Summary}
In ordinary language, the physical layer can be stated as follows:
\begin{quote}
Geometry determines which transitions are admissible. Quantum theory distributes the possible transitions. The primary interface creates the very possibility of distinction. Therefore an automatically structured interface field can self-organize these possibilities. Coherent fields form a manifold. The manifold stabilizes as a physical structure.
\end{quote}

\section{Geometric Admissibility}
Let \((\M,g)\) be spacetime with metric \(g\). The domain of causally admissible transitions is
\begin{equation}
\Adm_c(g):=\{(x,y)\in\M\times\M\mid y\in J_g^+(x)\}.
\end{equation}
The first layer is therefore
\begin{equation}
\boxed{g\Rightarrow\Adm_c(g).}
\end{equation}

\section{Quantum Distribution of Possibilities}
Over the domain of admissible transitions we construct the Hilbert space
\begin{equation}
\Hcal_I(g):=L^2(\Adm_c(g),\B_g,\mu_g).
\end{equation}
If \(\Psi_I\in\Hcal_I(g)\), then the state defines a distribution of possibilities:
\begin{equation}
\boxed{\Psi_I\Rightarrow\Dist_{\mathrm{poss}}(\Psi_I;g).}
\end{equation}

\section{Axiomatic Source of Automaticity}
From TMI-Core 0.2 we have
\begin{equation}
\boxed{N_0\xrightarrow{I_0}\Delta_0\to(A,\neg A).}
\end{equation}
\begin{equation}
\boxed{\Valid(I)\Rightarrow\Auto(I).}
\end{equation}
\begin{equation}
\boxed{\Auto(I)\equiv(\Input(I)\land\Adm(I)\land\Dir(I)\Rightarrow\AutoAct(I)).}
\end{equation}
If a physical field is an interface field, then
\begin{equation}
\boxed{\Xi_I\in\Ispace_{\mathrm{phys}}\Rightarrow\Auto(\Xi_I)\Rightarrow\XiAuto.}
\end{equation}

\section{Automatically Structured Interface Field}
\begin{definition}
An automatically structured interface field is an operator field
\begin{equation}
\XiAuto\in\Gamma(\End(\Hcal_I)),
\end{equation}
with an internal law of self-updating:
\begin{equation}
\partial_\tau\XiAuto=\mathcal A_I(\XiAuto,g,\Psi_I,\Adm_c(g),\Dist_{\mathrm{poss}}(\Psi_I;g)).
\end{equation}
\end{definition}
The operator \(\mathcal A_I\) is the internal auto-structuring rule. It does not represent an external controller, but a law internal to the interface-field dynamics.

\section{Internal Structuring Theorem}
\begin{theorem}
If \(\XiAuto\in\Gamma(\End(\Hcal_I))\), its dynamics is autonomous, closed in the admissible space, and there exists a stability functional \(\mathcal L_I\geq0\) such that
\begin{equation}
\frac{d}{d\tau}\mathcal L_I[\XiAuto]\leq0,
\end{equation}
then
\begin{equation}
\boxed{\XiAuto\Rightarrow\Struct_I.}
\end{equation}
\end{theorem}

\begin{proof}
The field \(\XiAuto\) acts on the possibility space \(\Hcal_I\). Autonomy excludes an external controlling operator, closure preserves admissibility, and the non-increase of \(\mathcal L_I\) defines motion toward a stable regime. This regime is denoted by \(\Struct_I:=\Attr(\mathcal A_I)\). Therefore \(\XiAuto\Rightarrow\Struct_I\).
\end{proof}

\section{Projections}
The full field admits the projection chain
\begin{equation}
\boxed{\XiAuto\to\Xi_I\to\chi_I.}
\end{equation}
The scalar projection can be written as
\begin{equation}
\boxed{\chi_I^2(x)=\frac{1}{N_I}\Tr_I\left[(\XiAuto(x))^\dagger\XiAuto(x)\right].}
\end{equation}
The stable interface structure is represented by
\begin{equation}
\boxed{\Struct_I=\mathcal S_I(g,\Psi_I,\XiAuto,\Adm_c(g),\Dist_{\mathrm{poss}}(\Psi_I;g)).}
\end{equation}

\section{Scaling}
A local automatically structured field scales into a network and then into a manifold:
\begin{equation}
\boxed{\XiAuto\to\{\Xi_{I,k}^{\mathrm{auto}}\}_{k\in K}\to\NXi\to\XiMan\to\Struct_I(\M).}
\end{equation}
The manifold of automatically structured interface fields is
\begin{equation}
\boxed{\XiMan=(V_\Xi,E_\Xi,O_\Xi,\mathcal A_\Xi,\PiStr,\PiAmp,\PiStruct).}
\end{equation}
Here \(V_\Xi\) is the set of local automatically structured fields, \(E_\Xi\) is the set of relations between them, \(O_\Xi\) is the order of coherence, \(\mathcal A_\Xi\) is the global law of auto-structuring, and \(\PiStr,\PiAmp,\PiStruct\) are structural, amplitude, and stability projections.

\section{Coherence and Synchronization}
For linked fields, define
\begin{equation}
\Coh_\Xi(k,l)=1-\Loss_\Xi(\Xi_{I,k}^{\mathrm{auto}},\Xi_{I,l}^{\mathrm{auto}}).
\end{equation}
The manifold condition is
\begin{equation}
\Coh_\Xi(\XiMan)\to1,\qquad \Sync_\Xi(\XiMan)\to1.
\end{equation}
This condition expresses that local fields should not merely coexist; they should become mutually coherent and synchronized enough to form a stable manifold.

\section{Computational Support and Engineering Projection}
The computational validation module tests an engineering projection of this architecture in quantum-control models:
\begin{equation}
I_Q^{coh-auto}\Rightarrow T_2^{eff}\uparrow\Rightarrow P_{success}\uparrow.
\end{equation}
The extended test line also introduces a safety boundary:
\begin{equation}
\Auto(I)\not\Rightarrow\operatorname{Safe}(I),
\end{equation}
with the engineering recovery principle
\begin{equation}
\Auto(I)+\operatorname{SafetyLayer}\Rightarrow\operatorname{SaferAdaptiveControl}.
\end{equation}
This computational support does not prove the existence of \(\XiAuto\) in nature. It supports the internal engineering idea that an automatic interface can improve observable quantum-control metrics under appropriate feedback and safety conditions.

\section{Status of Claims}
\begin{itemize}
\item Definitions: \(\Valid(I)\Rightarrow\Auto(I)\).
\item Internal consequences: \(\XiAuto\land\Stab(\XiAuto)\Rightarrow\Struct_I\).
\item Physical hypotheses: \(\XiAuto\) exists as a physical mechanism.
\item Experiments: \(T_2^{TMI}>T_2^{base}\) in a quantum implementation.
\end{itemize}

\section{Conclusion}
TMI-Physics 2.2 formulates the physical layer as a projection of the updated core of automatic interfaces. The scalar field \(\chi_I\) is preserved as the minimal observable projection, while the full form \(\XiAuto\) describes automatic structuring of distributed possibilities and scaling toward an interface manifold.

\begin{thebibliography}{9}
\bibitem{TMIcore} Salkutsan Aleksey. \textit{Theory of Manifold Interfaces}. \href{https://doi.org/10.5281/zenodo.19940675}{doi:10.5281/zenodo.19940675}.
\bibitem{TMIphysics21} Salkutsan Aleksey. \textit{Theory of Manifold Interfaces: A Regularized Effective Physical Layer, Observable Consequences, and Limitations}. 2026.
\bibitem{TMIvalidation} Salkutsan Aleksey. \textit{Computational Validation of Adaptive Interface Control in TMI Quantum Systems}. 2026.
\end{thebibliography}
\end{document}
